% problem-000059 1 铱氧簇异构与氧化还原
\section*{1. 铱氧簇异构与氧化还原}
\textit{下列为分问。}

\subsection*{1-1}
在该系列化合物中， $[ \mathrm { I r O } _ { 4 } ] ^ { + }$ 有 A、B、C 三种异构体。A ⽆对称中⼼，Ir—O 键的键⻓均为 170.8 pm。B中Ir的配位数为4，氧化态为+7，只有两种Ir—O键，键⻓分别为188.8 pm和168.0 pm。C中只有⼀个镜⾯，也只有两种 Ir—O 键，键⻓分别为 209.6 pm 和 167.9 pm。画出 A，B，C 的结构。

\subsection*{1-2}
$\mathrm { I r B r _ { 3 } }$ 与 $\mathrm { N } _ { 2 } \mathrm { O } _ { 5 }$ 可在室温下反应，得到⼀紫⾊配合物 $\mathrm { I r _ { 3 } O ( N O _ { 3 } ) _ { 1 0 } }$ ，同时产⽣⼀红棕⾊的液体单质，并释放出具有顺磁性的⽆⾊⽓体。写出该反应的⽅程式。

\subsection*{1-3}
具有三⽅双锥构型的 $\mathrm { C r } ( \mathrm { C O } ) _ { x } ( \mathrm { N O } ) _ { y }$ 配合物，满⾜18电⼦规则。计算�和�。
